% Part: first-order-logic
% Chapter: beyond
% Section: introduction

\documentclass[../../../include/open-logic-section]{subfiles}

\begin{document}

\olfileid[hi]{fol}{byd}{int}

\olsection{अवलोकन}

प्रथम-कोटि तर्कशास्त्र ही रुचिकर तार्किक तंत्र नहीं है: उसके अनेक
विस्तार और रूपांतर हैं। किसी तर्कशास्त्र में सामान्यतः एक भाषा का औपचारिक
विनिर्देशन, प्रायः---पर सदा नहीं---एक निगमन-तंत्र, और प्रायः---पर सदा
नहीं---एक अभीष्ट अर्थविज्ञान होता है। किंतु इस पद का तकनीकी प्रयोग एक
स्पष्ट प्रश्न उठाता है: जो तर्कशास्त्र प्रथम-कोटि तर्कशास्त्र नहीं हैं,
उनका सहज या दार्शनिक अर्थ में प्रयुक्त ``तर्क'' शब्द से क्या संबंध है?
नीचे वर्णित सभी तंत्र किसी-न-किसी प्रकार के विचार-विन्यास का प्रतिरूपण
करने के लिए बनाए गए हैं; क्या हम बता सकते हैं कि उन्हें तार्किक क्या बनाता
है?

इसका कोई सरल उत्तर नहीं मिलता। ``तर्क'' शब्द अलग-अलग ढंग से और अलग-अलग
संदर्भों में प्रयुक्त होता है; और ``सत्य'' की धारणा की तरह इस धारणा का भी
अनेक दार्शनिक दृष्टियों से विश्लेषण हुआ है। उदाहरण के लिए, तार्किक
विचार-विन्यास का लक्ष्य यह निर्धारित करना माना जा सकता है कि कौन-से कथन
अनिवार्यतः सत्य हैं, प्रागनुभविक रूप से सत्य हैं, अतार्किक पदों की व्याख्या
से स्वतंत्र सत्य हैं, अपने रूप के कारण सत्य हैं, अथवा भाषिक रूढ़ि के कारण
सत्य हैं; इनमें से प्रत्येक संकल्पना को काफी स्पष्ट करना पड़ता है। यदि हम
ध्यान केवल गणित में प्रयुक्त तर्कशास्त्र तक सीमित रखें, तब भी उसके विस्तार
पर बहुत कम सहमति है। उदाहरणतः \textit{Principia Mathematica} में रसेल और
व्हाइटहेड ने फ्रेगे द्वारा आरंभ की गई {\em तर्कवादी} परंपरा में तर्क के
आधार पर गणित विकसित करने का प्रयास किया। उनका तार्किक तंत्र नीचे वर्णित
उच्चतर-प्रकार तर्कशास्त्र से मिलता-जुलता था। अंततः उन्हें ऐसे अभिगृहीत
लाने पड़े जो अधिकांश मानकों पर शुद्धतः तार्किक नहीं लगते (विशेषतः अनंतता
का अभिगृहीत और अपचयनीयता का अभिगृहीत); फिर भी कोई यह मान सकता है कि
उच्चतर-कोटि विचार-विन्यास के कुछ रूपों को तार्किक स्वीकार करना चाहिए। इसके
विपरीत क्वाइन---जिनका सत्तामीमांसा-सिद्धांत ``प्रतिज्ञप्तियों'' को विचार के
वैध विषय नहीं मानता---तर्क देते हैं कि द्वितीय-कोटि और उच्चतर-कोटि
तर्कशास्त्र वास्तव में भेड़ की खाल में समुच्चय-सिद्धांत हैं; अर्थात् विधेयों
पर परिमाणीकरण वाले तंत्र शुद्धतः तार्किक नहीं हैं।

फिलहाल ऐसे दार्शनिक प्रश्नों को किसी और दिन के लिए छोड़ना और नीचे के
तंत्रों को तार्किक अथवा अन्य प्रकार के विचार-विन्यासों के औपचारिक आदर्शीकरण
मानना ही उचित है।

\end{document}
